\clearpage\thispagestyle{empty}
\pdfbookmark[0]{Quellen zu Gauß’ Briefen an Bode über Juno von 1804}{bodejuno1804-sources}
\begin{center}\textbf{Quellen zu Gauß’ Briefen an Bode über Juno von 1804}\end{center}
\noindent Diese KI-Übertragung in heutiges Deutsch folgt den Auszügen aus Carl Friedrich Gauß’ Briefen an Johann Elert Bode vom 30. September und 12. November 1804 in \href{https://gdz.sub.uni-goettingen.de/id/PPN236007467}{Gauß’ \textit{Werken}, Band~VI} (Göttingen, 1874), S.~248--250. Bode druckte die Auszüge unter dem Titel \textit{Beobachtung der Juno und Elemente ihrer Bahn, von Hr. Doct. Gauß in Braunschweig} im \textit{Astronomischen Jahrbuch für das Jahr 1808}, \href{https://www.digitale-sammlungen.de/en/view/bsb10538337?page=197}{S.~187}, \href{https://www.digitale-sammlungen.de/en/view/bsb10538337?page=198}{S.~188}, \href{https://www.digitale-sammlungen.de/en/view/bsb10538337?page=199}{S.~189} und \href{https://www.digitale-sammlungen.de/en/view/bsb10538337?page=200}{S.~190}. Das \href{https://www.digitale-sammlungen.de/en/view/bsb10538337?page=7}{Titelblatt} des Jahrbuchs für 1808 nennt Berlin, 1805, als Erscheinungsort und Erscheinungsjahr. Herausgeber war Johann Elert Bode.
\par\smallskip\noindent Übertragung und Satz wurden mit KI erstellt. Die Digitalisate des Jahrbuchs stammen von der Bayerischen Staatsbibliothek.
\clearpage
